\section{Nombres rationnels}
Les paramètres de procédures et les constantes que nous avons utilisés jusqu'à
présent ont des valeurs entières. En \python, on dit qu'ils sont de {\em type}\index{Type}
{\idt int} (de l'anglais \underline{int}eger, entier). Contrairement à beaucoup
d'autres langages, \python\ n'impose pas de limite à la taille des entiers~: {\idt 2 ** 100}
(2 à la puissance 100) vaut $1\,267\,650\,600\,228\,229\,401\,496\,703\,205\,376$, ce qui laisse
de la marge pour compter les côtés d'un polygone ou les carreaux d'une grille. Ils sont codés
en base deux dans la mémoire de l'ordinateur, sur autant de bits qu'il le faut.

Nous pouvons également utiliser en \python\ des nombres décimaux avec une virgule comme {\idt 3.14159} appelés flottants (de type \texttt{float}). Ils
s'écrivent selon la convention anglo-saxonne  avec un point à la place de notre
virgule. Les mathématiciens appellent ces nombres des {\em nombres rationnels}\index{Rationnel}. En \python, on
distingue les nombres entiers des nombres rationnels par la présence du point. Par exemple
{\idt 1000} est un entier alors que {\idt 1000.0} est un nombre rationnel, que l'on peut
aussi écrire {\idt 1000.} sans chiffre après le point. De même {\idt .5} est permis, mais
il est plus lisible d'écrire {\idt 0.5}.

Les nombres entiers et rationnels sont codés de manières différentes dans la mémoire
de l'ordinateur. Les nombres rationnels sont de type {\idt float}\footnote{\ De
l'anglais {\em floating point number}, nombre à virgule flottante. L'ordinateur
garde en mémoire un nombre fixe de chiffres significatifs et la position de la virgule
(qui ``flotte'' parmi ces chiffres). Ce ne sont pas des nombres
réels puisque les nombres de type {\idt float} ne peuvent pas avoir une infinité de chiffres
différents de zéro après la virgule$\:$!}.

\section{Paramètres de type ``float''}
En \python, on n'indique pas le type des paramètres d'une
procédure~: un même paramètre peut recevoir, d'un appel à l'autre, un entier ou un
rationnel, et c'est seulement pendant l'exécution, quand la valeur est utilisée,
que \python\ vérifie qu'elle convient (on dit que \python\ est un langage à
{\em typage dynamique}\index{Typage dynamique}). Par exemple, dans la procédure~:

\begin{center}
{\mcl def}\ {\idt arc(l):\ lg(l);\ vd()}
\end{center}

\noindent le paramètre formel {\idt l} (lire la lettre $\ell$ et non le chiffre $1$) est passé à la commande {\idt lg} du \robot\ qui
attend un nombre, entier ou rationnel.

Le paramètre effectif correspondant peut donc être indifféremment un entier, comme dans
{\idt arc(100)}, ou un nombre rationnel, écrit avec un point, comme dans {\idt arc(66.666666667)}~:
\python\ transforme de lui-même l'entier {\idt 100} en rationnel {\idt 100.0}
quand c'est nécessaire. Si l'on veut faire cette transformation soi-même, on utilise la
fonction {\idt float}\index{float}~: si {\idt n} est un entier,
{\idt float(n)} est le nombre rationnel correspondant (par exemple
{\idt float(3)} vaut {\idt 3.0}). C'est rarement utile.

Les procédures de commande du \robot\ {\idt lg}, {\idt lgX} et {\idt lgY} (ainsi que
{\idt lgRT}, {\idt rt} et {\idt tl} que nous verrons au chapitre suivant) acceptent des paramètres
rationnels. C'est pourquoi on peut écrire {\idt lg(20)} aussi bien que {\idt lg(20.5)}.
Les commandes {\idt avf}, {\idt cf}, {\idt ec}, {\idt deplacer} ou {\idt tracer}
ont au contraire des paramètres entiers~: {\idt avf(2.5)} arrête le programme avec le message
`{\tt TypeError:\ 'float' object cannot be interpreted as an integer}'
(un objet de type {\idt float} ne peut pas être interprété comme un entier).

\begin{exemple}[Spirale arithmétique]{\rm
La spirale ci-dessous est constituée d'arcs de cercles de rayons successifs
$100$$\:$; $100 / 1.5$ = $66.666\ldots$$\:$; $66.666\ldots / 1.5$ = $44.444\ldots$$\:$;
$44.444\ldots / 1.5$ = $29.629\ldots$$\:$; etc~:

 {\raggedright
 \hbox to 0pt{\hspace*{2cm}\vbox to 0pt{\includevspacePostScript[4cm]{SpiraleArithmetique}{0.8}{0.8}}}
 \input python/SpiraleArithmetique.tex
 }
\noindent $\Box$ }\end{exemple}

\section{Expressions rationnelles}\index{Expression rationnelle (réelle)}
L'addition, la soustraction, la multiplication et la division de nombres rationnels se
notent {\idt +}, {\idt -}, {\idt *} et {\idt /}, c'est-à-dire comme pour les entiers. Par exemple
{\idt (((2 * (10.3 + 0.7)) - 1.0) / 3)} vaut {\idt 7.0}

Quand une opération porte sur un entier et un rationnel, \python\ transforme l'entier en
rationnel et le résultat est un rationnel~: {\idt 2 * 3.5} vaut {\idt 7.0}.
Quand les deux opérandes sont entiers, {\idt +}, {\idt -} et {\idt *} donnent un entier.
La division, elle, demande de l'attention~: {\idt /} donne {\em toujours} un rationnel,
même entre deux entiers. Ainsi {\idt 7 / 2} vaut {\idt 3.5} et {\idt 6 / 2} vaut {\idt 3.0}.
La division entière, qui donne un entier, se note {\idt //}~: {\idt 7 // 2} vaut {\idt 3}.
Grâce à ces deux notations, le lecteur du programme sait toujours quelle division est faite.
Comme {\idt 6 / 2} est un rationnel, {\idt avf(6 / 2)} est refusé par la commande {\idt avf}
qui attend un entier~; il faut écrire {\idt avf(6 // 2)}.

Si $r$ est un nombre, {\idt sqrt($r$)} est sa racine carrée $\sqrt{r}$ (de l'anglais
\underline{sq}uare \underline{r}oo\underline{t}, racine carrée),
{\idt $r$ ** 2} est son carré $r^{2}$ (que l'on peut aussi écrire {\idt $r$ * $r$}),
{\idt abs($r$)} est sa valeur \underline{abs}olue (pour un entier comme pour un rationnel),
{\idt int($r$)} est sa partie entière obtenue en supprimant les chiffres après la virgule
(on dit qu'elle est tronquée) et {\idt round($r$)} est l'entier le plus proche (de
l'anglais \underline{round}, arrondir). Les résultats de {\idt int} et de {\idt round}
sont des entiers. Par exemple
{\idt int(0.3)} = {\idt int(0.5)} = {\idt int(0.9)} = {\idt 0} tandis que
{\idt round(0.3)} = {\idt 0} et {\idt round(0.5)} = {\idt round(0.9)} = {\idt 1}\footnote{\ Le
{\idt round} d'origine de \python\ arrondit les demis vers l'entier pair le plus proche
({\idt round(0.5)} vaut {\idt 0} et {\idt round(2.5)} vaut {\idt 2}). La bibliothèque du
\robot\ le remplace par un {\idt round} qui arrondit les demis à l'opposé de zéro, comme
Pascal et \ocaml~: {\idt round(2.5)} vaut {\idt 3} et {\idt round(-2.5)} vaut {\idt -3}.}.
La fonction {\idt sqrt} est fournie par la bibliothèque du \robot\ (ligne
{\mcl from}\ {\idt LibRobot}\ {\mcl import}\ {\idt *}), les autres font partie de \python.
Comme pour les procédures, l'argument d'une fonction est toujours placé entre
parenthèses~: {\idt sqrt(x + 1)}.

Dans une expression rationnelle, on peut utiliser la longueur horizontale {\idt valLgX()} et la
longueur verticale {\idt valLgY()} d'un carreau de la grille de déplacement du \robot. Par exemple
{\idt lgX(valLgX() * 2)} double la taille horizontale d'un carreau de la grille.

\begin{exemple}[Spirale d'or]{\rm
La spirale d'or ci-dessous est constituée d'arcs de cercles successifs dont les rayons
ont pour rapport le {\em nombre d'or\/}
$\Phi = \frac{1}{2}(1 + \sqrt{5}) \simeq 1.618033989$. Pour les tenants d'esthétique
traditionnelle, la spirale d'or constituerait une figure géométrique très harmonieuse
à l'œil.

{\raggedright
 \vbox to 0pt{\includevspacePostScript[4cm]{SpiraleDor}{0.95}{0.95}}
 \input python/SpiraleDor.tex
 }
\noindent $\Box$ }\end{exemple}

\section{Erreurs d'arrondi}
Pour les nombres réels\index{Réel}, qui ont un nombre infini de chiffres après la virgule, comme
$\frac{1}{3} = 0.3333333333333333\ldots$ ou $\pi = 3.1415926535897932\ldots$
l'ordinateur ne peut garder dans sa mémoire que quelques chiffres après la virgule, ce
qui conduit, comme pour une calculette, à des erreurs d'arrondi dans les calculs. Par
exemple en ma\-thé\-ma\-ti\-ques on a~:
\[\left(\left(\left(\left(\left(\sqrt{\sqrt{\sqrt{\sqrt{\sqrt{2}}}}}\right)^{2}\right)^{2}\right)^{2}\right)^{2}\right)^{2} =
\left(\left(\left(\left(\sqrt{\sqrt{\sqrt{\sqrt{2}}}}\right)^{2}\right)^{2}\right)^{2}\right)^{2} =
\] \[\left(\left(\left(\sqrt{\sqrt{\sqrt{2}}}\right)^{2}\right)^{2}\right)^{2} =
\left(\left(\sqrt{\sqrt{2}}\right)^{2}\right)^{2} =
\left(\sqrt{2}\right)^{2} =
2\]
tandis que l'ordinateur trouve  que~:
\begin{center}
{\idt ((((sqrt(sqrt(sqrt(sqrt(sqrt(2))))) ** 2)\hspace*{6em}}\\
{\idt \hspace*{6em}** 2) ** 2) ** 2) ** 2\ = 1.9999999999999964}
\end{center}

\noindent De même, il trouve que {\idt 0.1 + 0.2} vaut
{\idt 0.30000000000000004} (le programme {\idt ErreursDarrondi} du dossier
{\idt livre/\allowbreak 19\_expressions\_rationnelles} fait ces deux calculs).

\noindent  On retrouve des problèmes d'arrondi similaires avec la grille de dé\-pla\-ce\-ment du \robot. De
manière idéale, les coordonnées cartésiennes du \robot\ devraient être réelles.
L'ordinateur en conserve en mémoire une approximation rationnelle. Sur l'écran elles doivent
être arrondies à l'entier le plus proche pour correspondre à un pixel de l'écran.

\noindent \rule{137mm}{0.1mm} %---------------------------------------------
\begin{exercice}{Écrire des programmes \python\  pour réaliser les motifs de la figure
ci-dessous (seul le motif central qui utilise les fonc\-tions tri\-go\-no\-mé\-tri\-ques in\-tro\-dui\-tes
dans l'addenda du chapitre \ref{ch:rec}  est difficile).
\noindent $\Box$}\end{exercice}
\begin{figure}[htbp]
 \includevspacePostScript[4.6cm]{MotifsDeBeard}{0.95}{0.95}
 \caption{Motifs de Horemis et Beard$\:$; spirale de Baravelle}\label{fig:beard}
\end{figure}
\vfill
\begin{corrige}{
Nous donnons un programme, les autres se trouvant dans le dossier
{\idt livre/\allowbreak 19\_expressions\_rationnelles}~:
\label{prog:CarresImbriquesPeints}
\input python/CarresImbriquesPeints.tex
\noindent $\Box$}\end{corrige}

%%%%%%%%%%%%%%%%%%%%%%%%%%%%%%%%%%%%%%%%%%%%%%%%%%%%%%%%%%%%%%%%%%%%%%%%%%%
\noindent \framebox[\textwidth]{
\begin{minipage}{\myparboxwidth}
\setlength{\parindent}{2em}
\small \sf
\vspace{\abovedisplayskip}
\section*{Addenda}
Les nombres rationnels sont souvent qualifiés de réels par abus de langage. Ils
peuvent se noter en utilisant la {\em notation scientifique}\index{Notation scientifique}, qui
consiste à faire
suivre la partie fractionnaire par un facteur d'échelle formé de la lettre {\idt e} ou
{\idt E} (qui se lit ``dix à la puissance'') suivie d'un exposant.

\noindent Par exemple $-3.0${\idt e}$2$ = $-300.0$,
$3.0${\idt e}$0$ = $3.0$ et $3.0${\idt e}$-2$ = $0.03$. Un nombre rationnel ne contient pas de blanc. Par conséquent
$-3.0\hspace{0.3ex}\blanc\hspace{0.3ex}${\idt e}$\hspace{0.3ex}\blanc\hspace{0.3ex}2$ est erroné. En \python, les
nombres de type {\idt float} sont représentés avec 15 à 16 chiffres significatifs en base 10, le plus petit
rationnel strictement positif (sans perte de chiffres significatifs) étant de l'ordre de
$2.2${\idt e}$-308$ et le plus grand de l'ordre de $1.8${\idt e}$308$. Au-delà, le résultat
d'un calcul vaut {\idt inf} (l'infini)~: {\idt 1e308 * 10} vaut {\idt inf}.

\vspace{\abovedisplayskip}
\end{minipage}
}
%%%%%%%%%%%%%%%%%%%%%%%%%%%%%%%%%%%%%%%%%%%%%%%%%%%%%%%%%%%%%%%%%%%%%%%%%%%
%%%%%%%%%%%%%%%%%%%%%%%%%%%%%%%%%%%%%%%%%%%%%%%%%%%%%%%%%%%%%%%%%%%%%%%%%%%
\noindent \framebox[\textwidth]{
\begin{minipage}{\myparboxwidth}
\setlength{\parindent}{2em}
\small \sf
\vspace{\abovedisplayskip}
\smallskip

\index{Expression rationnelle (réelle)}\index{Expression entière}La valeur d'une expression arithmétique
peut être de type {\idt int} ou {\idt float}, selon les règles suivantes~:
\begin{list}{$\scriptscriptstyle\bullet$}{
}
\item les constantes entières (sans point) et les opérateurs $+$, $-$, $\ast$, {\idt //},
{\idt \%}, $\ast\ast$ (avec un exposant positif) et {\idt abs}\footnote{\ Nous utilisons des typographies différentes
pour les {\em mots-clés}\index{Identificateur réservé} (comme {\mcl def}
ou {\mcl for})
et les {\em identificateurs prédéfinis}\index{Identificateur prédéfini}
(comme {\idt sqrt} ou {\idt abs}). Les identificateurs introduits dans un programme \python\
(comme les noms de constantes, de procédures, etc.) doivent être
différents des mots-clés mais peuvent être identiques aux identificateurs prédéfinis. Cette
pratique est déconseillée car les programmes qui changent le sens habituel des identificateurs
prédéfinis sont illisibles.} appliqués à des entiers donnent un résultat entier
de type {\idt int}, de même que {\idt int} et {\idt round} qui
s'appliquent à un rationnel$\:$;
\item les constantes rationnelles (avec un point ou un exposant {\idt e} ou
{\idt E}), la division {\idt /} et la fonction {\idt sqrt} donnent un résultat rationnel
de type {\idt float}, de même que {\idt float} qui s'applique à un entier$\:$;
\item les opérateurs $+$, $-$, $\ast$, $\ast\ast$ et {\idt abs} appliqués à des rationnels,
ou à un entier et un rationnel, donnent un résultat rationnel.
\end{list}
Une expression qui mélange les deux, comme {\idt 1 + 2.5}, est donc permise~: l'entier est
transformé en rationnel et le résultat {\idt 3.5} est rationnel.

\smallskip

En \python, le type des paramètres n'est pas déclaré. \python\ transforme de lui-même
un entier en rationnel quand il le faut, mais jamais un rationnel en entier~: c'est au
programmeur de le demander avec {\idt int} ou {\idt round}, par exemple
{\idt avf(round(x))}. On peut indiquer le type attendu d'un paramètre formel lors de la
déclaration de la procédure, comme dans {\mcl def}\ {\idt arc(l:\ float):\ lg(l);\ vd()},
ce qui n'est pas nécessaire (\python\ ne le vérifie pas) mais peut aider le lecteur du programme.

\vspace{\abovedisplayskip}
\end{minipage}
}
%%%%%%%%%%%%%%%%%%%%%%%%%%%%%%%%%%%%%%%%%%%%%%%%%%%%%%%%%%%%%%%%%%%%%%%%%%%
